%% 300_ZweiRollen_Kfrak_De_ch.tex
%% FFGFT -- Dok. 300: Spiral-Time als Auslesung und entprellter Comparator;
%%                    der geteilte Korrekturfaktor K_frak = 1 - 100 xi (74/75)
%% Autor: Johann Pascher
%% Datum: Juli 2026
%% Sprache: Deutsch

\section*{Dok.~300 \quad Spiral-Time als Auslesung und entprellter Comparator
--- und der geteilte Korrekturfaktor $K_{\text{frak}} = 1-100\xi$ ($74/75$)}

\subsection*{Anlass und Abgrenzung}

Aus der Zusammenarbeit stammt die Beobachtung, Spiral-Time sei nicht ein
einzelner Operator, sondern eine Operator-Klasse mit mindestens zwei Rollen:
einer \emph{Auslesung} (was geschieht im Signal bzw.\ im Carrier) und einem
zweiten Teil, der entscheidet, ob diese Auslesung selbst als stabiler
Übergang gelten darf. Dieses Dokument präzisiert die zweite Rolle mess- und
schaltungstechnisch und ordnet den Korrekturfaktor, der beide Seiten
verbindet, seiner \emph{tatsächlichen} Herkunft im Korpus zu. Abgrenzung:
Kategorie-B-Befund (echte Kompatibilität / Kandidat-Überlapp), keine
HLV-interne Ableitung, keine physikalische Validierung.

\subsection*{Die zwei Rollen, präzise gefasst}

Die erste Rolle, die \emph{Auslesung}, ist ein Messwert: das Signal am
$+$Eingang eines Vergleichs. Die zweite Rolle ist kein Verwaltungsbegriff
(\enquote{Governance}), sondern präzise ein \textbf{Comparator mit Hysterese}
--- ein Schwellenentscheid, dessen Ausgang ein Bit trägt: gilt der Übergang
oder nicht. Zwei Formen der Hysterese:
\begin{itemize}
  \item \textbf{Amplituden-Hysterese}: zwei Schwellen im Pegel, ein Totband
    auf der Werteachse, das gegen Rauschen am Umschaltpunkt stabilisiert.
  \item \textbf{Zeit-Hysterese (Entprellung)}: der Ausgang akzeptiert eine
    Flanke erst, wenn sie über ein Zeitfenster \emph{stehen bleibt}. Das
    Totband liegt auf der Zeitachse.
\end{itemize}
Die zweite Rolle von Spiral-Time ist die Zeit-Hysterese. Das in der finiten
Proxy-Implementierung verwendete Fenster ($M=12$, $\tau=6$) ist ein
gewichteter Entpreller: ein Übergang gilt erst, wenn er das Fenster übersteht.
\enquote{Stabiler Übergang} ist damit die Zeit-Hysterese.

\subsection*{Gültigkeitsbedingung: Referenz aus Ableitung, nicht aus Kopie}

Ein Comparator trägt nur Information, wenn die Referenz am $-$Eingang aus einer
\emph{anderen} Quelle stammt als das Signal am $+$Eingang. Liegt an beiden
Eingängen dieselbe Größe, ist der Ausgang bedeutungslos. Genau dies war der
zirkuläre Kern des U3-Randfall-Gates in der ersten HLV-seitigen
Implementierung: die Referenz war eine Kopie des Auslese-Fensters
(\texttt{hard = hlv.copy()}), womit die Korrelation konstruktionsbedingt $1{,}0$
ist. Die FFGFT-seitige Prüfung (Dok.~299) ersetzt dies durch einen
\emph{abgeleiteten} Zeugen (BLP-Rückfluss $5{,}125$, Dok.~283), der vom Objekt
unabhängig und von einem Hard-Reset-Fenster nicht tragbar ist. Selbstreferenz
ist zulässig, wenn sie durch eine Ableitung geschlossen wird; zirkulär, wenn
durch eine Kopie.

\subsection*{Warum \enquote{dieselbe Spirale} kein Postulat ist}

Die Aussage über den Korrekturfaktor setzt voraus, dass das projizierte
HLV-Objekt und die FFGFT-Rekursion \emph{dieselbe} Spirale sind. Diese
Identität ist nicht angenommen, sondern durch mehrere unabhängige
Berührungspunkte gestützt:
\begin{itemize}
  \item \textbf{$2/9$}: dieselbe Umverteilung als Ausgabe der
    $Z_3$-Circulant-Diagonalisierung --- gefunden, nicht gesucht.
  \item \textbf{$D6$ / $\varphi$-Ikosaeder}: dieselbe Objektstruktur,
    $\tau=\varphi$, alle fünfzehn Achsenwinkel $\arccos(1/\sqrt5)=63{,}435^\circ$.
  \item \textbf{Zeit}: der Spiralzeit-Sektor als beschränkter Ausschnitt der
    aufgerollten Rekursionszeit (R47).
\end{itemize}
$2/9$ und $D6$ sind geometrisch, voneinander unabhängig und beide gefunden;
unter dieser Mehrfach-Verankerung ist \enquote{dieselbe Spirale} die einzige
verbleibende Lesart (Kategorie~B).

\subsection*{Das Schließungs-Komma und $K_{\text{frak}}$}

Die aufgerollte Rekursion schließt sich erst nach $N=1/(100\xi_0)=75$ gleichen
Teilstücken ($\xi=4/30000$), nicht nach einer Umdrehung (Dok.~295/297). Die
Näherung \enquote{eine Umdrehung deckt sich} vernachlässigt einen kleinen
Faktor: $75$ jeweils um $\varepsilon$ verkürzte Umdrehungen ergeben die exakte
Schließung, $75\,(1-\varepsilon)=74$, also
\begin{equation}
  \varepsilon \;=\; \tfrac{1}{75} \;=\; 100\,\xi \;=\; 0{,}0133\overline{3}
  \qquad(\text{als Winkel } 2\pi/75 = 4{,}8^\circ \text{ pro Umdrehung}).
\end{equation}
Der Pro-Schritt-Faktor der Rekursion ist damit
\begin{equation}
  K_{\text{frak}} \;=\; 1 - \varepsilon \;=\; 1 - 100\,\xi \;=\; \tfrac{74}{75}
  \;=\; 0{,}9867,
\end{equation}
in genau der Form, die der Korpus durchgängig führt (Dok.~133/231/032/285/018/005).
Die reziproke Streckenverhältnis-Form (verkürztes Stück $\to$ ganze Umdrehung)
ist $75/74 = 1{,}0135$.


\subsection*{Der tatsächliche Anker: nicht die $75$, sondern Dok.~133}

Die Herkunft von $100\xi$ --- und damit von $75$ \emph{und}
$K_{\text{frak}}$ --- liegt eine Ebene tiefer und ist unabhängig:
\begin{itemize}
  \item Die \textbf{Skalen-Rekursion} $\xi_{n+1}=\xi_n(1-100\xi_n)$
    (Dok.~133/275) hat den Pro-Schritt-Faktor $K_{\text{frak}}=1-100\xi$;
    $75=1/(100\xi_0)$ ist ihre Windungszahl.
  \item Der \textbf{Faktor $100$} emergiert aus dem RG-Lauf von
    $D_f=3-\xi$ über $\sim 17$ Größenordnungen (Planck $\to$ elektroschwach),
    effektiv $D_f^{\text{eff}}\approx 2{,}973$ (Dok.~133); geometrisch
    $100=4\cdot 5^2$ (Dok.~042).
  \item $K_{\text{frak}}$ ist \textbf{nicht willkürlich}, sondern durch die
    Forderung fixiert, dass zwei unabhängige Herleitungen des Massenverhältnisses
    $m_e/m_\mu$ denselben Wert liefern (Dok.~133, \emph{Eindeutige Bestimmung}).
  \item $\xi=4/30000$ selbst ist über Higgs-Relation, $\alpha$ und Koide
    verankert.
\end{itemize}
Das Schließungs-Komma \emph{re-exprimiert} damit die bereits verankerte
$100\xi$; es leitet sie nicht neu ab. Die $75$ ist ein Output derselben
Rekursion, kein unabhängiger Anker.

\subsection*{Konsequenz für die Projektion (Marcel)}

Damit bleibt die Aussage an Marcel gültig, nur richtig gebucht: \emph{wenn}
sein projiziertes, geschlossenes Umdrehungs-Objekt ein Ausschnitt derselben
Rekursion ist (belegt über $2/9$, $D6$, Zeit), dann trägt die Umrechnung
zwischen seiner Projektion und der aufgerollten exakten Kurve genau den
Pro-Schritt-Faktor $K_{\text{frak}}=1-100\xi=74/75$ (bzw.\ reziprok $75/74$)
--- einen Faktor, der \emph{unabhängig} verankert ist (Dok.~133), nicht aus der
$75$ und nicht frei gewählt. Es ist damit eine Konsistenz-Bedingung, keine
Ableitung aus der Windungszahl: verschluckt seine Projektion das Komma und
erscheint das Objekt sauber geschlossen, ist $K_{\text{frak}}$ die Umrechnung,
die es an der Grenze zur SI-Ebene wieder einträgt (Dok.~298, SM als
dekompaktifizierte Projektion). Die Einbettung $\iota$ (Dok.~299) ist an dieser
Stelle diese $K_{\text{frak}}$-Umrechnung.

\subsection*{Die Verzahnung, nicht der Ausschnitt}

Der Zusatz \enquote{ein Ausschnitt derselben Rekursion} ist bewusst unbestimmt.
Die $75$ Teilstücke sind unter der geteilten diskreten Symmetrie
gleichberechtigt: $75 = 3\cdot 5^2$ trägt genau die beiden Anker, die auch das
gemeinsame Objekt erzeugen --- die $Z_3$-Dreifachheit (dieselbe wie in
$2/9=2/3^2$) und die ikosaedrische Fünffachheit (dieselbe $5$ wie in
$\arccos(1/\sqrt5)$ und in $100=4\cdot5^2$, Dok.~042). Die Symmetrie bildet die
Teilstücke aufeinander ab; keines ist ausgezeichnet.

Damit ist die Frage, \emph{welches} der $75$ Teilstücke das projizierte Fenster
trifft, nicht offen, sondern gegenstandslos: eine ausgezeichnete Zuordnung
verbietet gerade die Symmetrie. Behauptet wird allein die \emph{Existenz der
Verzahnung} --- dass das geschlossene Objekt überhaupt als Teilstück dieser
Teilung sitzt, deren Winkelstruktur $D6$ festlegt und deren Zähler sich zu
$3\cdot5^2$ fügt ---, nicht eine Index-Zuordnung.

Für $K_{\text{frak}}$ ist das folgenlos und gerade der Grund, warum die
Umrechnung ohne Identifikation funktioniert: $K_{\text{frak}}=1-100\xi$ ist der
Pro-Schritt-Faktor der Teilung, über alle $75$ Schritte \emph{derselbe}, eben
weil sie gleichberechtigt sind. Er braucht kein ausgewähltes Teilstück, um
definiert zu sein; die bloße Existenz der Verzahnung genügt, damit die
Umrechnung zur aufgerollten Kurve $1-100\xi$ ist.

\subsection*{Grenze und Status}

Die Kette ist so belastbar wie ihre Verankerung: $2/9$ und $D6$ sind als
geteilte Objekte gebucht (Kategorie~B), $K_{\text{frak}}=1-100\xi$ ist über
Dok.~133 unabhängig verankert. Prüfskript des Kommas:
\texttt{Dok300\_Skripte/ffgft\_300\_kfrak\_schliessungskomma\_pruefung.py};
Zeitsektor-Gegenlauf: \texttt{Dok299\_Skripte/ffgft\_299\_iota\_gegenlauf\_winding\_locked.py}.
